% ============================================================
%  R. Nielsen, Religionsphilosophie (Kjøbenhavn: Gyldendal, 1869)
%  INDLEDNING  (§§ 1–3, pp. 1–30)
%  TRANSCRIPTION (Danish) — from KB scan, text layer extracted
% ============================================================
\documentclass[12pt,a4paper]{article}
\usepackage[utf8]{inputenc}
\usepackage[T1]{fontenc}
\usepackage[danish]{babel}
\usepackage{microtype}
\usepackage{setspace}
\usepackage[margin=1.2in]{geometry}
\usepackage{fancyhdr}
\usepackage{hyperref}

\hypersetup{
  pdftitle={Religionsphilosophie — Indledning},
  pdfauthor={R. Nielsen},
  hidelinks
}

\pagestyle{fancy}
\fancyhf{}
\rhead{Nielsen, \emph{Religionsphilosophie} (1869) — Indledning}
\cfoot{\thepage}

\setstretch{1.3}

\title{\textbf{Religionsphilosophie}\\[6pt]
  \large Indledning}
\author{R.\ Nielsen}
\date{Kjøbenhavn: Gyldendal, 1869, s.\ 1--30}

\begin{document}
\maketitle
\thispagestyle{fancy}

\bigskip

\begin{center}
\itshape
Tro og Viden ere som Principper absolut ueensartede.

\smallskip

Principernes Ueensartethed gjør intet Brud paa de
almeengyldige Tankelove, men viser Umuligheden af en
hvilkensomhelst videnskabelig Overgang enten fra Viden
til Tro eller fra Tro til Viden.

\smallskip

Den lange Strid mellem Religion og Videnskab har
oprindelig sin Grund i en uklar Blanding af de ueensartede
Principer, og maa altsaa høre op, naar Blandingen ophører.

\smallskip

Den Sætning, at Sandhed i Religionen ikke er Sandhed i
Videnskaben, er ikke længer et Paradox, naar Videns absolute
Grændse falder sammen med det Mysterium, hvorfra
Troeserkjendelsen gaaer ud, og hvortil den vender tilbage.

\smallskip

Det er disse Hovedsætninger, som Religionsphilosophien i
den foreliggende Form vil udvikle og begrunde. Om
Udviklingen og Begrundelsen har nogen Gyldighed, maa
Prøvelsen afgjøre.

\bigskip
Kjøbenhavn, den 28de December 1868.

R.\ Nielsen.
\end{center}

\bigskip\bigskip

\section*{INDLEDNING}

Religionsphilosophiens Forhold til Religionen kan ikke paa
tilfredsstillende Maade angives i almindelige
Begrebsbestemmelser; dets eiendommelige Natur udkræver en
eiendommelig Behandling. Hvad der siges religiøst: Viis mig
din Tro af dine Gjerninger, det maa siges philosophisk: Viis
mig dit Princip af dets Gjennemførelse.

\subsection*{Religion og Philosophie}
\subsubsection*{§ 1}

Er Philosophien i det Hele en Videnskab, saa maa
Religionsphilosophien, saasandt den skal fortjene Navn af
Philosophie, naturligviis ogsaa være en Videnskab, men just
det Særegne, at den er en Videnskab om Religionen, giver den
et fra de øvrige Videnskaber afvigende Præg og paatrykker
den et eget Mærke. Hvad dette betyder, vil bedst kunne
forstaaes ved et Blik paa de Hovedformer,
Religionsphilosophien i nyere Tid har antaget.

Den videnskabelige Form maa bestemmes ved Principet.
Grundspørgsmaalet er: have Religion og Philosophie det samme
Princip, eller ere Principerne forskjellige? Dersom Religionens
Princip er forskjelligt fra Philosophiens, hvilket Princip skal
Videnskaben da følge, Philosophiens eller Religionens? Først
ved en Besvarelse af disse Spørgsmaal kan
Religionsphilosophiens Begreb og Methode bestemmes.

Gaaer man ud fra den Forudsætning, at Religion og Philosophie
have samme Princip, saa bliver Religionsphilosophien en
speculativ Videnskab; antager man, at Principerne vel ere
forskjellige, men at Sandheden dog maa bedømmes efter
Vidensprincipet, saa bliver Religionsphilosophien en kritisk
Videnskab. Antager man endelig Principerne for absolut
ueensartede og lægger det religiøse Princip til Grund ved den
videnskabelige Belysning, saa bliver Religionsphilosophien en
omvendt Videnskab. Den første af disse Synsmaader
repræsenteres af Hegel; den anden fornemmelig af
L.\ Feuerbach; den tredie er den, som i nærværende
Fremstilling kommer til Anvendelse.

\medskip

\noindent\textit{a) Religionsphilosophien som speculativ Videnskab.}

\medskip

Religionens saavelsom Philosophiens Gjenstand er, ifølge
Hegel, den evige Sandhed i dens Objectivitet, Gud og Intet
uden Gud og Explicationen af Gud. Philosophien udvikler kun
sig selv, idet den udvikler Religionen, og idet den udvikler
Religionen, udvikler den sig selv. Saaledes falder Religion og
Philosophie da sammen til Eet. Philosophien er i Virkelighed
Gudstjeneste, er Religion, thi den er Forsagelse af subjective
Indfald og Meninger i Beskjæftigelsen med Gud. Ikke desto
mindre er Beskjæftigelsen med Gud i Philosophien en anden,
foregaaer paa en anden Art og Maade end i Religionen. I
Religionen er Gud Bevidsthedens umiddelbare Gjenstand; i
Philosophien er Bevidstheden om Gud Guds egen Bevidsthed om
sig selv: Forskjellen ligger i Formen. Religionen som Religion
er folkelig, er for Alle og just derfor henviist til
Forestillingens billedlige Form; men en saadan Form er
uadæqvat for Indholdet. Formen er endelig, Indholdet
uendeligt: dette er et Misforhold. Kan dette Misforhold hæves;
kan det samme Indhold, uden at forandre sig, fremstilles i to
aldeles forskjellige Former?

Her er Vanskeligheden, og paa denne Vanskelighed strander
Hegel. Thi hvor det gjælder Sandheden som Sandhed, gaaer
Ideen saaledes op i Formen, at Formen umulig kan forandres,
uden at Indholdet forandres. Under sin Tankebevægelse bliver
den speculative Religionsphilosophie umærkelig ført til et
andet Resultat end det, hvorpaa den oprindelig sigtede.
Bevægelsen er denne: det begynder med, at Videnskaben skal
anerkjende og begrunde Religionens Sandhed; Religionens
Sandhed bestemmes ved det Sande i Indholdet, d.\ e.\ ved det
større eller mindre Quantum Livsidealitet, en given Religion i
Forhold til det Trin af Formudvikling, hvorpaa den befinder
sig, er istand til at udtrykke; da nu Religionens Form ikke er
videnskabelig, men den videnskabelige Form den eneste
fuldkomne, saa bringes det til Bevidsthed, at selv den
fuldkomne Religion har en ufuldkommen Form, og kan paa Grund
af Formens Ufuldkommenhed kun udtrykke det Sande paa
ufuldkommen Maade. Det, hvorved den fuldkomne Religion
adskiller sig fra Philosophien, er altsaa en Form, der maa
ophæves; men idet Formen ophæves, ophæves det for Religionen
Eiendommelige og dermed Religionen selv. Ved at miste sin Form
mister Religionen tillige sit Indhold, thi Indholdets Sandhed
beroer paa Formen. Altsaa ender det med, at Philosophien, for
at kunne begribe det Sande i Religionen, maa begribe Religionen
selv som en Usandhed. Det negative, fra Videnskabens Side
frastødende Forhold mellem Philosophie og Religion er med
kraftig Conseqvens kommen til Udtalelse i Nyhegelianismen.

\medskip

\noindent\textit{b) Religionsphilosophien som kritisk Videnskab.}

\medskip

Medens den speculativt positive Retning nærmest holder sig til
det for Religion og Philosophie fælles Indhold og stræber at
udvikle dette Sandhedens sande Indhold til den sande Form,
gaaer den negativt kritiske Retning, i Erkjendelse af
Religionens og Philosophiens principielle Forskjellighed,
nærmest ud fra den for Religionen eiendommelige Form, paaviser
dens øiensynlige Strid med Fornuften, og kommer saa gjennem en
Paaviisning af det Usande i Formen til det Usande i Indholdet.

Umiddelbart seer det ud, som havde Religionen med lutter
virkelige, om end overnaturlige, Kjendsgjerninger at gjøre; men
den til Viden udviklede Fornuft begriber, at de overnaturlige
Kjendsgjerninger ikke ere virkelige, og de virkelige ikke
overnaturlige. Det er ikke Aanden som Aand, men det sjælelige
Livs ideale Naturkræfter, der røre sig oprindeligt i den
oprindeligt religiøse Bevidsthed. De religiøse Forestillinger
ere vel ikke Frembringelser af bevidst Vilkaarlighed, tvertimod
de skyldes en ubevidst Digtning, en Sjælenes Trang til
Selvobjectivering, en psychologisk Nødvendighed, hvormed den af
Ideen overvældede Phantasie har stillet det Guddommelige i
Mennesket frem for Mennesket ved at stille det udenfor
Mennesket. Religionen er menneskelig; den er den menneskelige
Fornufts første elementære Gjennembrud; den har sit Udspring,
ikke i Guds Kjærlighed, men i Menneskets Lidenskaber.

Hjertets Drømme, Sindets Rørelser, Sjælens Længsler indeholde
de Grundelementer, af hvilke den ideebeaandede Phantasie har
dannet alle Troens Guder, ogsaa den ene Almægtige. Dette
stadfæstes ved Indblik saavel i Monotheismens som i
Polytheismens Væsen. Den mægtige, vrede, hævnfnysende Jehova i
det Gamle Testamente er, videnskabelig seet, en idealiseret
Speiling af Jødefolkets abstracte, mægtige Aand, af Israels
stærke, lovbundne og dog selvraadige Jeg. Christus i det Nye
Testamente er ikke et virkeligt Gudmenneske, thi et virkeligt
Gudmenneske er en Modsigelse, men dog væsentlig en Menneskegud,
det saarede Hjertes, den ængstede Sjæls forsonende Gud, en Gud
for de Lidende. For Phantasie og Følelse fremstiller det Ideale
sig umiddelbart, derpaa beroer ikke blot Ønsket selv, men ogsaa
dets Opfyldelse. Mennesket med sin indsnevrende
Naturbegrændsning ønsker sig Magt; havde det Magt som det har
Villie, vilde det føle sig saligt og frit; men det er afmægtigt;
saa føler det i sin Afmagt Trang til guddommelig Bistand. Den
guddommelige Bistand bestemmes umiddelbart saaledes, at en Gud,
der har Magt over Naturen, vilkaarlig griber ind i Tingenes Gang
og hjælper ved et Mirakel. I Miraklet anskues den ubegrændsede,
fra Naturtvang frigjorte Villie, en Villie, der formaaer at
tilintetgjøre al Modstand, tilfredsstille alle Ønsker og høre
alle Bønner. Miraklet er Troens Grund, Haabets Udvei,
Phantasiens og det fromme Gemyts Trøst, hellige Eventyr; naar
en Gud er bleven saa gammel og svag, at han ikke længere kan
gjøre Mirakler, gjør han bedst i at resignere, ophøre at være
Gud, og ligesaa maa en Religion, der ikke vedblivende kan
opvise Mirakler, ophøre at være Religion. Her er det
Vendepunkt, hvor Striden mellem Religion og Philosophie
culminerer. Religionen holder paa Miraklet, Philosophien paa
Fornuften. Miraklet er i Modsigelse med Fornuften; naar
Fornuften vaagner til Erkjendelse af Virkelighedens Love,
forsvinder Miraklet, og med Miraklet tillige Religionen.

Saaledes har da Philosophien, som negativ Kritik, ifølge egen
Forsikkring, fuldbyrdet Religionens Opløsning.

Ad den negative Vei er der imidlertid naaet et positivt
Resultat. Resultatet er, at Religionen, forskudt og forhaanet
af »Videnskaben«, bliver henviist til sit eget Princip og
derved opfordret til at stole paa egne Kræfter og hævde sin
Selvstændighed overfor Videnskaben. Det afgjørende Brud mellem
Religion og Videnskab kan da paa disse Vilkaar ligesaavel komme
Religionen som Videnskaben tilgode. Religionens absolute
Ueensartethed med al objectiv Viden og Videnskab er med
gjennemtrængende Skarphed og seirrig Overlegenhed paaviist af
S.\ Kierkegaard.

Grundspørgsmaalet, hvorfra S.\ Kierkegaard gaaer ud, er dette:
Hvorvidt kan Sandheden læres? Her møder saa, hvad Sokrates i
\textit{Meno} kalder en stridslysten Sætning: at et Menneske
umuligt kan søge, hvad han veed, og ligesaa umuligt kan søge,
hvad han ikke veed, thi hvad han veed, kan han ikke søge, da
han veed det, og hvad han ikke veed, kan han ikke søge, thi han
veed jo end ikke, hvad det er, han skal søge. Sokrates hæver
Vanskeligheden derved, at al Læren og Søgen kun er en Erindren,
saa at den Uvidende blot behøver at mindes, for ved sig selv at
besinde sig paa, hvad han veed. Sandheden bringes saaledes ikke
ind i ham, men var i ham. Dette er, hvad vi kunne kalde den
humane Forklaring. Fra dette Synspunkt faaer »Øieblikket«, det
Guddommeliges Indtræden i Tiden, Underet, ingen Betydning. »For
den sokratiske Betragtning er ethvert Menneske sig selv det
Centrale, og hele Verden centraliserer kun til ham, fordi hans
Selverkjendelse er en Gudserkjendelse.« Den, der lærer
Sandheden, lærer den ved sig selv, Læreren er kun Anledningen;
Øieblikket, da Sandheden indlyser, er i Forhold til den
værende, Menneskets Væsen iboende Sandhed et forsvindende og
forsaavidt ligegyldigt Tidspunkt. »Udgangspunktet er et Intet;
thi i samme Øieblik som jeg opdager, at jeg fra Evighed af har
vidst Sandheden uden at vide det, i samme Nu er hiint Øieblik
skjult i det Evige, indoptaget deri saaledes, at jeg, saa at
sige, end ikke kan finde det, selv om jeg søgte det, fordi her
intet Her og Der er, men kun et \textit{ubique et nusquam}.«

Anderledes fra det exclusivt religiøse Standpunkt. Her
forudsættes Mennesket at være, ikke i Sandheden, men i
Usandheden. Skal nu den, der ved Slægtens og ved egen Skyld har
tabt Sandheden, komme til Sandheds Erkjendelse, nemlig den
Sandheds Erkjendelse, der fører til Salighed, da bliver det
nødvendigt, at Gud ikke alene oplyser den Formørkede om
Usandheden, men tillige ved selv at indtræde i Tiden og
existere som Menneske giver det i Usandheden værende Menneske
Betingelsen for at naae Sandheden. Ved Guds Selvmeddelelse i
Tiden faaer Øieblikket, det bestemte Tidspunkt, og dermed
Underet, absolut afgjørende Betydning. Underet er Forstandens
Paradox. Ved at fastholde Underet og opstille Paradoxet har
S.\ Kierkegaard sikkret Religionen mod Videnskabens Overgreb og
reist et Bolværk mod den negative Kritiks Indbrud; men saa er
han ogsaa med afgjørende Sigte paa sit »Indledningsproblem, ikke
til Christendommen, men til det at blive en Christen«, bleven
staaende ved den Afgjørelse, at Christendommen ikke er nogen
Lære, men en Existens-Meddelelse, og en Existens-Meddelelse,
udtrykt i det Problem: »at Individets evige Salighed afgjøres i
Tiden ved Forholdet til noget Historisk, der ydermere er
saaledes historisk, at i dets Sammensætning Det er medoptaget,
hvilket ifølge sit Væsen ikke kan blive historisk, og altsaa maa
blive det i Kraft af det Absurde«.

Ifølge den negative Kritik er det Videnskaben, der erklærer
Religionen for et Absurdum og derved mener at opløse den;
ifølge den Kierkegaardske Opgjørelse er det Religionen, der,
den objective Viden til Trods, erklærer sig selv for et Absurdum
og derved een Gang for alle udelukker Videnskaben. Heraf sees,
at Religion og Philosophie ere ikke blot ueensartede, men
absolut ueensartede. Er nu med denne Forudsætning en
Religionsphilosophie mulig?

\medskip

\noindent\textit{c) Religionsphilosophien som omvendt Videnskab.}

\medskip

Det exclusive, ved Underet bestemte Gudsforhold, der er
Religionens egen uafviselige Forudsætning, kan Videnskaben ikke
tilegne sig. Denne Forudsætning er det Ubegribelige, det
Uudgrundelige, og kan kun anerkjendes paa den Betingelse, at
den menneskelige Viden anerkjender en Skranke, en absolut
Grændse, hvorved den maa standse. Hvad der i andre
Erkjendelsens Egne viser sig modsigende og urimeligt, er altid
et Lavere, som Videnskaben frit kan hæve sig over; den kan
drøfte dette Lavere, opløse det, og ved dets Opløsning bestandig
skride videre frem. Men Religionen, det for Viden Modsigende, er
ikke et Lavere, men et Høiere, thi Religionen er Livets høieste
ideale Magt; hvor denne Magt i Medfør af sit oprindelige Væsen
møder med »et Absurdum«, der er Veien objectivt ufremkommelig,
der finder Fornuften Gjennemgangen spærret, der maa Videnskaben
bøie af.

Men det for den videnskabelige Tænkning Absurde, »Paradoxet«,
falder dog indenfor den tænkende Bevidsthed, saasandt det falder
indenfor Religionen; hvoraf igjen følger, at det for Videnskaben
Utænkelige kan være tænkeligt i Religionen. Der er en væsentlig
Forskjel mellem det: at overskride Videns Grændse og at
overskride Tankens Grændse, mellem det: at falde udenfor Viden
og at falde udenfor Tanken. Forskjellen mellem Tænkeligt og
Utænkeligt bliver her at bestemme ved det Princip, hvorfra
Tænkningen gaaer ud. At Paradoxet falder udenfor Viden, er
vist; men falder det derfor udenfor Tanken? For at Tanken kan
støde sig paa Paradoxet, maa Tænkningen jo selv bringe Paradoxet
frem; Paradoxet er kun utænkeligt derved, at det bliver tænkt;
dets dialektiske Knude er en Sammenslyngning af absolut
ueensartede Tanketraade. Existentielt er Knuden uløselig; men
deraf følger ingenlunde, at den ogsaa skulde være det
intellectuelt. Den intellectuelle Løsning er betinget af, at de
ueensartede Traade opredes og skilles fra hverandre.
Videnskaben vilde ikke være istand til at opfatte og bestemme
sit Vidensprincip, dersom den ikke ved Modsætningen mellem det,
som kan vides, og det, som ikke kan vides, var istand til, om
end blot negativt, at charakterisere det for Viden
uigjennemtrængelige Troesprincip. Men er det Videnskaben muligt
at adskille de absolut ueensartede Principer, da er det den jo
ogsaa muligt at adskille de fra disse Principer udgaaende
Tanker, muligt altsaa at gjøre Forskjel mellem Videnstanker og
Troestanker, mellem Vidensbegreber og Troesbegreber.

Med den eiendommelige Opgave at klare Brydningsforholdet mellem
de ueensartede Tanker er Religionsphilosophien nærmest at
bestemme som en Grændsevidenskab.

Religionsphilosophien bliver imidlertid ikke staaende ved
almindelige Grændsebestemmelser for Tro og Viden; den kan
umulig lade Religionen udvikle sit Indhold uden at lade et heelt
System af Troesbegreber komme til Udvikling. Men Troens
religiøse Indhold er, ifølge Principet, et overvidenskabeligt
Indhold; her indtræder da den Modsigelse, at et
overvidenskabeligt Indhold optages i Videnskaben. Dersom
Videnskabsformen nu virkelig blev forvexlet med Indholdets sande
Form og indført og fastholdt istedenfor denne, vilde Modsigelsen
ganske vist være uopløselig. Men dette er ingenlunde Tilfældet.
Den videnskabelige Reflexion har just den Bøielighed, at den kan
bringe Vidensformen til at reflectere en anden og modsat Form.
Thi ligesom den klare Dag kan bringe Øiet til at glemme
Lysskjæret og see paa Gjenstandene, saaledes kan ogsaa
Reflexionen bringe Tanken til i Viden at see paa Troen og
glemme Viden. Religionsphilosophien er saa langt fra at
paatvinge det religiøse Indhold en fremmed Form, at dens
Formudvikling tvertimod bestaaer i at frigjøre det fra de
Tilsætninger af umiddelbart individuelle Forestillinger,
Billeder, Følelser, Villiesyttringer o.\ s.\ v., hvormed det i
Troeslivet selv altid er sammenvoxet. Forsaavidt Reflexionen,
saavel hvad Formen som hvad Indholdet angaaer, ved Viden
ophæver Viden, idet den lægger an paa at klare, ikke en
videnskabelig, men en overvidenskabelig Sandhed, er Forholdet
mellem Sandheden og Videnskaben væsentlig vendt om, og
Religionsphilosophien bliver med Hensyn hertil at charakterisere
som en omvendt Videnskab.

Det for enhver Videnskab Afgjørende er Begrundelsen. Nu kan jo
vistnok en grundig Indholdsudvikling betragtes som en
eiendommelig Sagudvikling, altsaa som en Udvikling, hvorunder
Sagen udfolder sine Momenter og begrunder sig selv. Men i alle
directe Videnskaber er Sagbegrundelsens Princip i Grunden Eet
med det almene Vidensprincip, saa at Sagbegrundelsen Punkt for
Punkt falder sammen med den paagjældende Videnskabs egen
Selvbegrundelse. I den omvendte Videnskab er Begrundelsen
derimod omvendt. Grundprincipet er her saa langt fra at gaae op
i Vidensprincipet, at de to absolut ueensartede Principer,
Underets og Fornuftens, overalt frastøde hinanden og i deres
negative Forhold bestemme hinanden negativt. Om dette negative
Forhold med de dertil svarende Vexelbestemmelser kan der nu
vistnok gives en klar Viden; men den almene, de modsatte
Principer oversvævende Viden er paa nærværende Trin ikke
længere istand til at tage Parti for sit eget Princip; den
omvendte Videnskab har just den Opgave, at sætte Troen i Kraft
og bringe Viden selv til at bøie sig for den. Som det gaaer med
det begrundende Princip, maa det gaae med Begrundelsen. Den
omvendte Videnskab skal hverken fortie eller fornegte de
almeengyldige Vidensgrunde, men den skal bringe til klar
Indsigt, hvorledes Religionens Selvbekræftelse netop foregaaer
derved, at Vidensgrundene conseqvent opløse sig og give Plads
for Troens absolut selvgyldige Grunde. Ved at begrunde
Vidensgrundenes Omsætning og Opløsning i Troesgrunde begrunder
Religionsphilosophien sin egen videnskabelige Opløsning og faaer
paa disse Vilkaar Betydning som en paa Religionen sig ofrende,
sig selv i Troeslære ophævende Videnskab.

\subsection*{Religion og Mythologie}
\subsubsection*{§ 2}

Af det Forhold, hvori Philosophien sætter sig til Religionen,
afhænger den videnskabelige Opfattelse af Modsætningen mellem
sand og falsk Religion, mellem Religion og Mythologie. Den
speculative Fornuft gjør Mythologien til Religion; den kritiske
Forstand gjør Religionen til Mythologie; kun i den omvendte
Viden kan den principielle Adskillelse mellem Religion og
Mythologie gjennemføres. Vi skulle betragte disse Synsmaader
nærmere.

\medskip

\noindent\textit{a) Mythologie er væsentlig Religion: den speculative Opfattelse.}

\medskip

Med den Forudsætning, at Sandheden i Religionen og Sandheden i
Videnskaben ikke ere to indbyrdes ueensartede Sandheder, men
den ene og samme absolute Sandhed, maa den speculative
Religionsphilosophie conseqvent anerkjende Religionen i enhver
Mythologie, og betragte den mythiske Indklædning som et Hylster,
der let kan borttages. Religionen er Aand, og for Aanden som
Aand ere de endelige Forestillinger ligegyldige. At udgrunde
Aandens Væsen er at udgrunde Religionens Væsen; for at udgrunde
Religionens Væsen udvikler Philosophien Religionens Begreb.
Religionens Begreb er i dets Almindelighed et Slægtsbegreb, men
efter dets særlige Fremtrædelsesmaader et Artsbegreb.
Religionens Idee har sondret sit Indhold og udlagt sine Momenter
i en Flerhed af mere og mindre eensidige Religioner, for
igjennem disse at samle sig ud af Adspredelsen og komme til
Fuldendelse i den fuldkomne, ene, sande Religion. »Det har«,
siger Hegel, »været Aandens Arbeide gjennem Aartusinder, at
udføre Religionens Begreb og gjøre den til Gjenstand for
Bevidstheden.« Opgaven er at vise, hvorledes Menneske-Aanden,
begyndende i al Naturalighed, har Trin for Trin kunnet hæve sig
over den naturbundne Tilstand til stedse større Frihed og
omsider naae det Vendepunkt, da Aanden begreb sig som Aand og
udtrykte Begrebet i den forbilledlige Existens, der under
Samfundets Form skal tilegnes og virkeliggjøres af Alle.

Hegels religionsphilosophiske Systematik er i denne Henseende
charakteristisk. Fra Naturreligionen, det Trin, hvor
Bevidstheden i gjærende Eenhed af Natur og Aand anskuer det
Guddommelige, gjennem den aandige Individualitets eller den
frie Subjectivitets --- Høihedens, Skjønhedens og
Hensigtsmæssighedens --- Religioner skrider Udviklingen frem til
Opfattelsen og Bestemmelsen af den i sig klare, aabenbare, til
Aabenbaringsbegrebet svarende Religion, den absolute Religion,
Christendommen.

Ifølge denne Synsmaade ere alle Religioner beslægtede, alle
udsprungne af een Rod, alle Grene paa een Stamme. Den gamle
Adskillelse imellem Naturreligionerne og den aabenbarede
Religion er hos Hegel overseet, endog i den Grad, at
»Høihedens Religion« i Jødedommen bliver for Begrebets Skyld
stillet lavere end »Skjønhedsreligionen« hos Grækerne og
»Hensigtsmæssighedsreligionen« hos Romerne. Hvad følger nu
heraf? At den jødiske Religion er Mythologie ligesaavel som den
græske og romerske! Men naar Jødedommen er Mythologie, da er
Christendommen ogsaa Mythologie: den absolute Religion er af
samme Oprindelse som de relative. Med den absolute Religion er
det Mythiske altsaa ikke forsvundet; det er tvertimod blevet
Aanden gjennemsigtigt som et til Religionens eiendommelige Form
Medhørende. Først naar det i Religionen forestillede Begreb
begribes i Viden, d.\ v.\ s.\ naar Religionen forvandles til
Philosophie, forsvinder det Mythiske.

»Den absolute Religion er den aabenbare. Religionen er da først
det Aabenbare, naar Religionens Begreb selv er for sig, d.\ v.\ s.\
naar dens Begreb er blevet sig selv objectivt, ikke i
indskrænket, endelig Objectivitet, men saaledes at den efter sit
Begreb er sig objectiv. Vi have saaledes Tvende: Bevidstheden og
Objectet; men i den med sig opfyldte Religion, den aabenbare,
der har fattet sig selv, er Indholdet selv Gjenstanden, og denne
Gjenstand, det sig vidende Væsen, er Aanden.«\footnote{Hegel,
\textit{Religionsphilosophie} II. Berlin 1840, s.\ 192--93.}

\medskip

\noindent\textit{b) Religion er Mythologie: den kritiske Opfattelse.}

\medskip

Striden mellem Religion og Philosophie begyndte hos Grækerne
meget tidlig, og hvad Romerne angaaer, kunde jo to Augurer paa
Ciceros Tid ikke betragte hinanden uden at lee. Vistnok afgav
Religionen det oprindelige Grundlag for Aandslivet hos Grækerne
saavelsom hos Romerne, men hos dem begge viser det sig ogsaa,
at de religiøse Forestillingers Form kun paa Folkelivets første,
barnlige Trin tilfredsstiller Aanden: Udviklingsloven er, at
Bevidstheden med den tiltagende Reflexion og den fremskridende
Cultur voxer fra Religionen. Da Sokrates havde opdaget
Subjectivitetens Dialektik, sat Selverkjendelsen som Philosophiens
Maal og forsaavidt indført »nye Guder«, kunde Folkereligionen
ikke længere modstaae Philosophernes Angreb. Blev end
Mythebegrebet selv ikke i Oldtiden udviklet med den Klarhed og
Dybde, som skeet er i Nutiden, saa blev det dog, hvad Gudernes
Liv og Heroernes Bedrifter angaaer, snart tilstrækkelig
indlysende, at »det Guddommelige«, for at bruge et Udtryk af
Strauss, »ikke kunde være skeet saaledes, at det saaledes Skete
ikke kunde være guddommeligt«. Denne Opløsning, der forklarer
Mytherne, men tilintetgjør Guderne, er ikke et aandeligt
Tilbageskridt, men i Tankens, Fornuftens, Selvbevidsthedens Navn
et væsentligt Fremskridt. Det Mythologiske er det ubevidst
Illusoriske; men dette ligefrem Illusoriske kan Religionen ikke
afføre sig uden at blotte sin medfødte Svaghed og give sig selv
til Priis. De overnaturlige Phænomener i det Nye Testamente ere
ligesaa fornuftstridige som Begivenhederne i den græske
Gudekrig. Religion og Mythologie ere eenstydige Begreber, og
Forskjellen mellem lavere og høiere, ufuldkomnere og fuldkomnere
Religioner kun en Forskjel mellem lavere og høiere, ufuldkomnere
og fuldkomnere Mythologier.

Jo høiere en Mythologie staaer, jo mere opfyldt den er af Væsen
og Aand, desto længere vil den naturligviis holde ud med
Culturbevidstheden; men just den Omstændighed, at den
aandfuldeste af alle Myther, Mythen om den Opstandne, har kunnet
bære Culturen gjennem atten Aarhundreder og holde ud med
Menneske-Aanden, indtil Bevidstheden omsider fik Bugt med alle
uklare Forudsætninger og kom til sig selv, idet den »som det
frie, for sig værende Almene, blev sig selv gjennemsigtig«,
netop dette viser, hvad det egentlig er, der i den
»aabenbarede« Religion bliver det Aabenbare. Det er ikke en
Grundforskjel mellem Mythe og Aabenbaring, der kommer for
Dagen; nei, det, der nu kommer for Dagen, er noget ganske Andet,
det er den store objective Sandhed, at Menneske-Aanden, ikke det
enkelte Menneskes Aand, men Menneske-Aanden i sit Princip, er
sin egen Guddom, er i den religiøse Selvfordobling baade Subject
og Object, som endelig Aand Subject, som Aand i indre
Uendelighed Object. Dette er ingen blasphemisk Selvforgudelse,
det er jo Opdagelsen af det Absolute, af Ideen som Idee, af det
Subjectives og det Objectives gudmenneskelige Eenhed.

Ved at gjennemføre Vidensprincipet har den philosophiske Kritik
conseqvent opløst Religionen som Religion --- thi at ville
bevare det Religiøse, naar Religionens Princip er opløst, er
naturligviis Nonsens --- den har da med det Samme naaet det
Grændseskjel, hvor Livsmagterne reagere: det Religiøse gjør sin
Ret gjældende, Troen træder op imod Viden, og de ueensartede
Principer maale sig med hinanden.

\medskip

\noindent\textit{c) Religion er absolut ueensartet med Mythologie: den antirationelle Opfattelse.}

\medskip

Der er saavel i den kritiske som i den speculative Philosophies
Stilling til Religionen en Tvetydighed, der fra Først til Sidst
maa blende Betragteren og forvirre hans Opfattelse. Denne
Tvetydighed ligger i Brugen af Ordet Sandhed. At den evige,
absolute Sandhed kun er een, og at den ene guddommelige Sandhed
ikke kan være bleven splidagtig med sig selv eller have deelt
sig i to med hinanden indbyrdes stridende Sandheder, er ganske
vist en ligesaa utvetydig som ubestridelig Forudsætning. Men
idet man saa uden videre overseer Forskjellen imellem en
guddommelig og en menneskelig Viden af Sandheden, overseer, at
hvad der i Guds Viden er absolut, maa i Menneskets være relativt,
overseer den principielle Forskjel mellem religiøs og
videnskabelig Erkjendelse, fremkommer Tvetydigheden. Det seer
ud, som om Philosophien virkelig gik ind paa Religionens egne
Forudsætninger, lod disse Forudsætninger udvikle deres religiøse
Conseqvenser, og derved kom til stridige, hinanden ophævende
Resultater. Ja, hvis det virkelig gik saaledes til, da kunde det
med Sandhed siges, at den philosophiske Kritik lod Religionen
selv fuldbyrde sin egen Opløsning. Men saaledes gaaer det ikke
til. Istedenfor at begynde med det religiøse Princip, begynder
jo Videnskaben tvertimod med Vidensprincipet. Ved at anvendes
paa Religionen bliver Vidensprincipet conseqvent et
religionsfornegtende Princip. Naar det da faaer Udseende af, at
Kritiken beviser Religionens Selvopløsning, saa er dette kun et
Skin. Sandheden er, at Kritiken faaer ud, hvad den selv har sat
ind; den begynder med at sætte Religionsfornegtelsen ind, og
ender med at faae Religionsfornegtelsen ud: den opløsende Kritik
bevæger sig i en Cirkel. Thi hvad er det vel, som Videnskaben i
Principet sætter som det Mythiske? Det er just det for
Religionen eiendommelige Gudsforhold. At klare og bestemme
Gudsforholdet er at klare og bestemme de for Aandslivet
afgjørende Grundbetingelser. Nu bør dog Videnskaben vide, at
disse Grundbetingelser stille sig anderledes paa Troens end paa
Videns Omraade. Er det virkelig Videnskabens Alvor at faae
Aandslivet selv fra Grunden af opklaret og forstaaet, da maa
den, idetmindste indtil videre, kunne vise den Selvfornegtelse,
at see bort fra sit eget Princip, og lade Religionen komme til
Orde.

At Religionsprincipet er absolut ueensartet med Videnskabens
Princip, kan sees af Begrundelsen. Alle Vidensgrunde ere i
dybeste Forstand Nødvendighedsgrunde, Religionens Grunde
derimod ere Villiesgrunde. Guds Villie er Religionens Absolute.
Skulde det, hvad Tankeløse paastaae, være en Modsigelse, at
Videnskaben bliver vidende om et med Viden absolut ueensartet
Princip, da maatte det netop være den negative Kritik, der har
bildt sig ind at have løst Modsigelsen. Hvad vil denne Kritik?
Opløse den positive Religion i Kraft af Viden. Hvorfor vil den
opløse den positive Religion? Fordi dens oprindelige
Forudsætninger ere absolut ueensartede med Videnskab.

\end{document}
